\providecommand{\Needspace}[1]{}
\Needspace{8\baselineskip}
\item \textbf{Modelo de decaimiento radiactivo con NumPy y Matplotlib}.
\nopagebreak[4]

\begin{tcolorbox}[colback=blue!2!white,colframe=blue!55!black,title={Enunciado},
                  before skip=3pt,after skip=5pt,boxsep=3pt,top=3pt,bottom=3pt,
                  breakable]
\small
\raggedright
El decaimiento radiactivo de una muestra puede modelarse mediante la ley
exponencial
\[
N(t)=N_0e^{-\lambda t},
\]
donde \(N(t)\) representa el número de núcleos radiactivos presentes en el
instante \(t\), \(N_0\) corresponde al número inicial de núcleos y \(\lambda\)
es la constante de decaimiento radiactivo.

Use los valores
\[
N_0 = 1000,
\qquad
\lambda = 0.23.
\]

Realice las siguientes operaciones:

\begin{itemize}
\setlength{\topsep}{4pt}
\setlength{\itemsep}{3pt}
\setlength{\parsep}{0pt}
\item[(a)] Importe \texttt{numpy} y \texttt{matplotlib.pyplot}. Cree un arreglo
\texttt{t} con 300 valores igualmente espaciados entre \(0\) y \(20\) usando
\texttt{np.linspace}.
\item[(b)] Cree una función llamada \texttt{decaimiento(t, N0, lam)} que retorne
\(N(t)\).
\item[(c)] Evalúe la función para obtener el arreglo \texttt{N}.
\item[(d)] Calcule el número mínimo de núcleos radiactivos, el número máximo y
el número promedio del arreglo \texttt{N}.
\item[(e)] Determine el tiempo aproximado en que \(N(t)\) es menor que
\texttt{250} por primera vez. Puede usar \texttt{np.where}.
\item[(f)] Grafique \(N(t)\) en función del tiempo usando \texttt{plt.plot}.
Agregue una línea horizontal en \(N=250\), etiquetas, título, leyenda y grilla.
\end{itemize}
\end{tcolorbox}

\begin{solution}
\begin{pythonjupyter}
# (a)
import numpy as np
import matplotlib.pyplot as plt

N0 = 1000
lam = 0.23

t = np.linspace(0, 20, 300)

# (b)
def decaimiento(t, N0, lam):
    N = N0 * np.exp(-lam * t)
    return N

# (c)
N = decaimiento(t, N0, lam)

# (d)
N_min = np.min(N)
N_max = np.max(N)
N_promedio = np.mean(N)

# (e)
indices_bajo_250 = np.where(N < 250)[0]
primer_indice = indices_bajo_250[0]
tiempo_bajo_250 = t[primer_indice]

print("Número mínimo:", N_min)
print("Número máximo:", N_max)
print("Número promedio:", N_promedio)
print("Primer tiempo con N < 250:", tiempo_bajo_250)

# (f)
plt.figure(figsize=(7, 5))
plt.plot(t, N, label="N(t)")
plt.axhline(250, color="gray", linestyle="--", label="N = 250")
plt.xlabel("Tiempo t")
plt.ylabel("Número de núcleos N(t)")
plt.title("Modelo de decaimiento radiactivo")
plt.legend()
plt.grid(True, alpha=0.3)
plt.show()
\end{pythonjupyter}
\end{solution}
